\documentclass[10pt,a4paper]{amsart}
\usepackage[margin=19mm]{geometry}
\usepackage{amsmath,amssymb,mathtools}
\usepackage[scr=rsfs,cal=euler]{mathalfa}
\usepackage{xfrac}
\usepackage[unicode,hidelinks,bookmarks=false]{hyperref}
\newcommand{\ie}{i.e.\ }
\newcommand{\cf}{cf.\ }
\newcommand{\resp}{respectively\ }
\newcommand{\Subsection}{\subsection}
\providecommand{\nobreakdash}{\nobreak\hbox{-}}
\setlength{\emergencystretch}{2em}
\multlinegap0pt
\usepackage{amssymb}
\usepackage{mathtools}
\usepackage{xcolor}
\usepackage{algorithm}\usepackage{algpseudocode}
\usepackage[shortlabels]{enumitem}
\newtheorem{proposition}{Proposition}
\usepackage{cleveref}
\crefname{proposition}{Proposition}{Propositions}
\newcommand{\R}{\mathbb{R}}
\newcommand{\ip}[2]{\langle #1, #2 \rangle}
\renewcommand{\vec}[1]{{#1}}
\title{Data-driven discovery of polynomial ODEs with provably bounded solutions}
\author{Albert Alcalde, Giovanni Fantuzzi}
\date{Source: arXiv:2604.26933v1 (CC BY 4.0)}
\begin{document}
\maketitle

\noindent\emph{Source and licence.} Data-driven discovery of polynomial ODEs with provably bounded solutions. Version arXiv:2604.26933v1 (CC BY 4.0), distributed by arXiv under the Creative Commons Attribution 4.0 International licence, http://creativecommons.org/licenses/by/4.0/, as recorded in the arXiv licence metadata for this exact version and shown on the abstract page https://arxiv.org/abs/2604.26933v1. Full licence notice: \texttt{licenses/arxiv-2604.26933-CC-BY-4.0.txt}.\par\smallskip

\noindent\textit{Editorial scope.} Counted unit: the article's proposition proposition (source0.tex lines 636--638). The article's complete original proof of it is delivered with it (source0.tex lines 639--674, one \texttt{proof} environment of 36 lines). No mathematics is written, repaired or re-derived: the statement and the proof are the article's own lines, in the article's own notation, and no retained line is substituted or re-flowed.  Retained uncounted as the source-local context the proof needs: lines 169--171; lines 378--384.  Nothing was omitted from any retained span, so the package carries no omission record.  No retained line contains a citation, so the wrapper carries no bibliography and no bibliography entry.  No retained line contains a figure environment, an image-inclusion command or drawing code, so no image is delivered and the package declares no asset dependency.  Editorial formatting only, in addition: an AMS \texttt{amsart} preamble that restates the article's own declarations and macros used by the retained lines, namely the environments \texttt{proposition} on separate counters, together with the standard packages those lines need.  The retained environments print in this document as Proposition 1 in order of appearance, whereas the article prints the same items with the numbering of its own full document.  The complete native source of the article as distributed in the arXiv e-print for this exact version is bundled unchanged as \texttt{sources/199436/source0.tex}.
\begin{equation}\label{e:lyapunov-inequality}
    b - v(\vec{x}) - \alpha\,  \ip{\vec{f}(\vec{x})}{\nabla v(\vec{x})}  \geq 0 \quad \forall \vec{x} \in \R^n,
\end{equation}

\begin{equation}\label{eq:ex3}
f^*(\vec{x}) =
\begin{pmatrix}
x_1 - x_2 - x_1 x_2^{2} + x_2^{3} \\
x_1 + x_2 - x_1^{2} x_2 - x_1^{3}
\end{pmatrix}.
\end{equation}

\begin{proposition}
    If $f$ is the vector field in \cref{eq:ex3}, there exist no $\alpha>0$, $b\in \R$, and coercive quadratic polynomial $v$ satisfying the Lyapunov inequality \cref{e:lyapunov-inequality}.
\end{proposition} 

\begin{proof}
    For the sake of contradiction, suppose that $v$ is a quadratic polynomial such that the Lyapunov inequality \cref{e:lyapunov-inequality} holds for some $\alpha>0$ and $b\in\R$.
    Since the vector field $\vec{f}$ in \cref{eq:ex3} is equivariant under the symmetry transformation $x\mapsto -x$, a standard symmetry reduction argument shows that $\frac12 v(\vec{x})+ \frac12v(-\vec{x})$ also satisfies \cref{e:lyapunov-inequality} for the same values of $\alpha$ and $b$. We may therefore assume without any loss of generality that $v(-\vec{x})=v(\vec{x})$, so it takes the form
    $v(x) = A x_1^2 + B x_1 x_2 + C x_2^2 + D$ 
    for some coefficients $A$, $B$, $C$ and $D$. We must also have $A,C>0$ because $v$ is coercive by assumption.
    
    Upon substituting this quadratic polynomial $v$ and the vector field $\vec{f}$ from \cref{eq:ex3} into the Lyapunov inequality \cref{e:lyapunov-inequality} we find that
    \begin{align*}
        0\leq 
        b - D
        - (A + B\alpha +2A\alpha)x_1^2 
        - (B+2B\alpha-2A\alpha+2C\alpha)x_1 x_2
        \\
        - (C - B\alpha +2C\alpha) x_2^2 
        + B\alpha x_1^4 
        + (B + 2C)\alpha x_1^3 x_2
        \\
        + 2(A+C)\alpha x_1^2 x_2^2
        + (B - 2A)\alpha x_1 x_2^3
        - B \alpha x_2^4
    \end{align*}
    for all $\vec{x}=(x_1,x_2)\in\R^2$.
    This means the coefficients of the monomials $x_1^4$ and $x_2^4$ must be nonnegative, so $B=0$. We then conclude that
    \begin{align*}
        0\leq 
        b - D
        - A(1 +2\alpha)x_1^2 
        - 2(C-A)\alpha x_1 x_2
        - C(1 +2\alpha) x_2^2 \\
        + 2C\alpha x_1^3 x_2
        + 2(A+C)\alpha x_1^2 x_2^2
        - 2A\alpha x_1 x_2^3
    \end{align*}
    for all $\vec{x}=(x_1,x_2)\in\R^2$. Since $C\alpha>0$, however,
    this inequality cannot be true because the polynomial on its right-hand side is cubic in $x_1$ for fixed $x_2\neq 0$, so it is not sign-definite.
\end{proof}

\end{document}
